% Fragmento sugerido para documentar la prueba v0.5.
% Actualizar los valores con los resultados reales.

\subsection{Resultados de prueba de carga moderada v0.5}

La quinta ejecución automatizada se enfocó en evaluar el comportamiento del backend
ante concurrencia moderada sobre endpoints funcionales no destructivos. La prueba no
ejecutó exportaciones pesadas; se limitó a registro, autenticación, consulta de sesión,
carga de página, consulta del área de estudio y consulta de estado de exportación.

\begin{table}[H]
\centering
\begin{tabular}{|p{4.5cm}|p{4cm}|p{4cm}|}
\hline
\textbf{Métrica} & \textbf{Resultado} & \textbf{Criterio} \\
\hline
Usuarios concurrentes & Pendiente & 10 \\
\hline
Solicitudes por usuario & Pendiente & 20 \\
\hline
Total de solicitudes funcionales & Pendiente & 200 \\
\hline
Tasa de éxito & Pendiente & $\geq$ 95\% \\
\hline
Solicitudes fallidas & Pendiente & 0 \\
\hline
Latencia p95 & Pendiente & $\leq$ 2000 ms \\
\hline
Resultado general & Pendiente & Aprobado/Fallido \\
\hline
\end{tabular}
\caption{Resumen de prueba de carga moderada v0.5}
\end{table}

Los endpoints evaluados fueron \texttt{/}, \texttt{/api/session}, \texttt{/api/study-area}
y \texttt{/api/area/geotiff-status}, usando sesiones autenticadas por cada usuario virtual.
La evidencia se almacena en \texttt{tests\_evidence/load/}.
